\section{Bloc Positions and Key Stakeholders}

\lead{It would be a mistake to treat NATO as a single voice on this issue. The Alliance agrees that the cuts are a threat, but it splits along a predictable line: the states most exposed to the cables want rapid, forceful response, while those with larger merchant fleets and more distant coastlines weigh the precedent any forceful action would set.}

Complicating this further, NATO itself is not legally the first responder. As officers inside the Alliance's own Baltic task force have stressed, the response to a suspected attack rests with the affected coastal state and its legal system. Cable owners hold primary responsibility for their security, with NATO there to assist rather than to provide blanket protection.

\subsection{Baltic States}

The frontline Baltic states, Estonia, Latvia, and Lithuania, are the most hawkish bloc. They are small, almost entirely dependent on a handful of cables for power and data, and they read every incident through the lens of Russian hybrid pressure on NATO's eastern flank. All three remain firmly committed to NATO and active in the Baltic Sentry.\footnote{Jackson School of International Studies. (2025). Baltic Sea undersea cable security. University of Washington}

After the Estlink 2 cut, Estonia's foreign minister framed the damage as an attack on vital infrastructure rather than an accident.\footnote{Finland boards oil tanker suspected of causing internet, power cable damage in Baltic. (2024, December 26). RGuťGrs.} Their problem is capability, not will. These states historically built land forces and left the sea to others, and their navies are limited, a gap exposed in May 2025 when the Estonian Navy approached the sanctioned tanker Jaguar, whose Gabon registration had just been revoked, but found how little margin a small navy has when a vessel refuses to comply and a Russian jet shadows the operation.\footnote{Estonia says Russia sent jet after attempt to stop sanction-breaking tanker. (2025, May 15). RGuťGrs}

Lithuania, which lost roughly a fifth of its internet capacity in the November 2024 cuts, has leaned on public-private cooperation, signing an agreement between its armed forces and the grid operator Litgrid to protect critical subsea links.\footnote{Lithuanian army, power grid operator to boost security of subsea infrastructure. (2025). BBC KGws}

\subsection{Finland and Sweden}

Finland and Sweden, the newest members, sit slightly apart as the enforcement test cases. Finland set the benchmark the rest of the Alliance now cites, escorting the Eagle S into its territorial waters, boarding it, and detaining ship and cargo, with President Stubb arguing that three cable incidents in a single year could not be coincidence.\footnote{Finland seizes Russian oil tanker after suspected undersea fiber-optic cable sabotage. (2024, December 27). USKI KGws. \url{https://news.usni.org/2024/12/27/finland-seizes-russian-oil-tanker-after-suspected-under} sea-fiber-optic-cable-sabotage}

Yet both countries also illustrate the limits of going hard. Sweden seized the Vezhen after the Gotland cable damage but released it once prosecutors concluded the anchor had likely dropped accidentally in high winds.\footnote{Sweden says ship broke Baltic Sea cable by accident. (2025, February 3). RGuťGrs. \url{https://www.reuters.com/world/europe/sweden-rules-out-sabotage-baltic-sea-cable-dam} age-case-2025-02-03/} A Finnish court later dismissed the case against the Eagle S crew, ruling that prosecutors had not proven intent and that any negligence had to be pursued by the flag state or the crew's home countries. The lesson for delegates is that even the toughest responders keep running into the same evidentiary wall.\footnote{Finland's Eagle S decision: How the deep sea goes dark when international law fails to protect states from underwater sabotage. (2025, December). Washingťon InťGrnaťional Law Journal. \url{https://wilj.org/2025/12/04/finlands-eagle-s-decisionwhen-the-deep-sea-goes-dark-whe} n-international-law-fails-to-protect-states-from-underwater-sabotage/}

\subsection{Larger Western Members}

The larger Western members carry the operational weight but move more cautiously.

\begin{itemize}
  \item Germany has become the regional coordinator, running Commander Task Force Baltic out of Rostock on behalf of NATO's Maritime Command, with Defence Minister Pistorius among the most willing to name Russia openly.\footnote{Braw, E. (2025). How the Baltic Sea nations have tackled suspicious cable cuts. Atlantic Council.}
  \item The United Kingdom leads the parallel Joint Expeditionary Force track, activating the AI-driven Nordic Warden system, shadowing Russian vessels such as the spy ship Yantar, and sanctioning more than a hundred shadow fleet ships in May 2025. \footnote{British-led force to use AI in tracking Russia's shadow fleet. (2025, January 7). DGfiGnsG KGws}
  \item The United States engaged early through Nordic-Baltic cyber and infrastructure coordination, though the depth of sustained American commitment is an open question delegates should probe.\footnote{Jackson School of International Studies. (2025). Baltic Sea undersea cable security. University of Washington}
  \item Denmark, holding the straits that are the only exits from the sea, is a natural gatekeeper but also belongs to the bloc of major maritime trading nations wary that aggressive interdiction in the Baltic could be turned against their own vessels elsewhere.\footnote{Braw, E. (2025). The shadow fleet is undermining the maritime order more brazenly than ever. Atlantic Council.}
\end{itemize}

\subsection{Russia}

Russia is the actor every position paper must address, and its stance is flat denial. The Kremlin has dismissed seizures as a matter of little concern and rejects responsibility for the incidents.\footnote{Finland says vessel suspected of cutting cable may be part of Russia's shadow fleet. (2024, December 26). ThG KGw York TimGs. \url{https://www.nytimes.com/2024/12/26/world/europe/finland-estonia-cables-russia.html}}

Behind the denial sits the shadow fleet, the instrument that makes the strategy work. Since 2022 Russia has assembled more than 400 ageing tankers used to export oil and, allegedly, to sabotage undersea cables.\footnote{Russia's shadow fleet: A masterclass in sanctions evasion. (2024). GGopoliťical Moniťor. \url{https://www.geopoliticalmonitor.com/russias-shadow-fleet-a-masterclass-in-sanctions-ev} asion/} More than 60\% of Russian seaborne crude passes through the Baltic. These vessels change flags often and carry no Western insurance, giving Moscow plausible deniability. Russia has also begun escorting these vessels and signalling a readiness to intervene militarily on their behalf, a dangerous step in an already securitised sea.\footnote{Estonia says Russia sent jet after attempt to stop sanction-breaking tanker. (2025, May 15). RGuťGrs}

\subsection{China}

China is the quieter but unavoidable stakeholder, present through the vessels themselves. Its pattern has been to deny, delay, then concede minimally: Beijing first denied the Newnew Polar Bear's involvement and only later admitted the ship caused the damage, attributing it to bad weather and an accident.\footnote{Beijing's Baltic confession exposes undersea vulnerability. (2024). ThG InťGrprGťGr, Lowy Institute. \url{https://www.lowyinstitute.org/the-interpreter/beijing-s-baltic-confession-exposes-unders} ea-vulnerability} Over the Yi Peng 3, China permitted a limited joint inspection after a month-long standoff while refusing to let the ship enter Swedish jurisdiction.\footnote{China lets Sweden, Finland, Germany and Denmark board ship in cable probe. (2024, December 19). RGuťGrs}

European concern about Chinese leverage over critical infrastructure runs wider than these incidents, with the European Parliament flagging the risks of Chinese firms acquiring European internet companies and of suppliers such as HMN Technologies.\footnote{A Chinese-flagged ship cut Baltic Sea internet cables. This time, Europe was watching. (2024, December). Carnegie Endowment for International Peace. \url{https://carnegieendowment.org/emissary/2024/12/baltic-sea-internet-cable-cut-europe-n} ato-security}

\subsubsection{Non-state actors and private operators}

The non-state layer is central rather than peripheral. The cables and pipelines are largely privately owned, by operators such as Cinia, Arelion, Litgrid, and Elisa, which means commercial firms sit on the front line of a security problem and bear first responsibility for repair and resilience.\footnote{Braw, E. (2025). How the Baltic Sea nations have tackled suspicious cable cuts. Atlantic Council}

The shadow fleet's enablers, the flag-of-convenience registries, opaque shell owners like the United Arab Emirates company behind the Eagle S, and the insurers and tracking systems these vessels exploit, are exactly where enforcement pressure is now aimed.\footnote{Russia-linked cable-cutting tanker seized by Finland "was loaded with spying equipment." (2025). Lloyd's Lisť. \url{https://www.lloydslist.com/LL1151955/Russia-linked-cable-cutting-tanker-seized-by-Finlan} d-was-loaded-with-spying-equipment}

The European Union has stepped in as an institutional actor in its own right, publishing an Action Plan on Cable Security in February 2025 and committing close to a billion euros to surveillance and a fleet of emergency repair vessels.\footnote{European Commission \& High Representative. (2025, February 21). EU Acťion Plan on CablG SGcuriťy (JOIN/2025/9). \url{https://eur-lex.europa.eu/legal-content/EN/TXT/?uri=celex:52025JC0009}} For delegates, the practical question is how a NATO committee divides labour with the EU, with private operators, and with the flag states whose cooperation any lasting solution requires.
